%杨舒云的实验报告编辑界面，使用了Huanyu Shi，2019级的模板，杨舒云在此拜谢ORZ!

%!TEX program = xelatex
\documentclass[dvipsnames, svgnames,a4paper,11pt]{article}
\input{YsySettings} 
\usepackage{lipsum}
\usepackage{adjustbox}

\begin{document}
	
	
	
	%实验报告封面	
	\begin{table}
		\renewcommand\arraystretch{1.7}
		\begin{tabularx}{\textwidth}{
				|X|X|X|X
				|X|X|X|X|}
			\hline
			\multicolumn{2}{|c|}{预习报告}&\multicolumn{2}{|c|}{实验记录}&\multicolumn{2}{|c|}{分析讨论}&\multicolumn{2}{|c|}{总成绩}\\
			\hline
			\LARGE20 & & \LARGE30 & & \LARGE50 & & \LARGE100 & \\
			\hline
		\end{tabularx}
	\end{table}
	
	\begin{table}
		\renewcommand\arraystretch{1.7}
		\begin{tabularx}{\textwidth}{|X|X|X|X|}
			\hline
			年级、专业： & 2022级 物理学 &组号： & 2\\
			\hline
			姓名： & 杨舒云  & 学号： & 22344020\\
			\hline
			实验时间： & 2023/11/16 & 教师签名： & \\
			\hline
		\end{tabularx}
	\end{table}
	
	\begin{center}
		\LARGE AA4 \quad 传感器综合实验
	\end{center}
	
	\textbf{【实验报告注意事项】}
	\begin{enumerate}
		\item 实验报告由三部分组成：
		\begin{enumerate}
			\item 预习报告：课前认真研读实验讲义，弄清实验原理；实验所需的仪器设备、用具及其使用、完成课前预习思考题；了解实验需要测量的物理量，并根据要求提前准备实验记录表格（可以参考实验报告模板，可以打印）。\textcolor{red}{\textbf{（20分）}}
			\item 实验记录：认真、客观记录实验条件、实验过程中的现象以及数据。实验记录请用珠笔或者钢笔书写并签名（\textcolor{red}{\textbf{用铅笔记录的被认为无效}}）。\textcolor{red}{\textbf{保持原始记录，包括写错删除部分，如因误记需要修改记录，必须按规范修改。}}（不得输入电脑打印，但可扫描手记后打印扫描件）；离开前请实验教师检查记录并签名。\textcolor{red}{\textbf{（30分）}}
			\item 数据处理及分析讨论：处理实验原始数据（学习仪器使用类型的实验除外），对数据的可靠性和合理性进行分析；按规范呈现数据和结果（图、表），包括数据、图表按顺序编号及其引用；分析物理现象（含回答实验思考题，写出问题思考过程，必要时按规范引用数据）；最后得出结论。\textcolor{red}{\textbf{（30分）}}
		\end{enumerate}
		\textbf{实验报告就是将预习报告、实验记录、和数据处理与分析合起来，加上本页封面。\textcolor{red}{（80分）}}
		\item 每次完成实验后的一周内交\textbf{实验报告}（特殊情况不能超过两周）。
	\end{enumerate}
	
	\textbf{【实验安全注意事项】}	
	\begin{enumerate}
		\item 所有实验均在 ELVIS 板上接线，调节、测量读数由电脑控制完成。
		\item \textcolor{red}{\textbf{每做一个小实验，都请先连好线路、并确保检查无误（电源正负极没有接反；电线都插紧；没有短路），
		再打开 ELVIS 板右上方 PROTO 电源，以防烧坏仪器。}}
		\item 每做完一个实验，请关掉右上 PROTO 电源、拆掉导线、将电路开关断路，再开始下一个实验连线。
		\item 电压/电阻/电流表请选择与待测信号最接近的量程，否则会有较大误差（仪器制造问题）。
		\item ELVIS 板四孔脚标 12 内部连通，脚标 34 内部连通，脚标 12 与脚标 34 之间不连通。
		
	\end{enumerate}
	
%	\textbf{【特别鸣谢及模板说明】}	
	
%	感谢2019级学长石寰宇为本实验报告提供\LaTeX 模板。
	
	
	
	\clearpage
	\tableofcontents
	\clearpage
	
	
	
	%预习报告	
	\setcounter{section}{0}
	\section{AA4 传感器综合实验 \quad\heiti 预习报告}
	
	\subsection{实验目的}
	\begin{enumerate}
		\item 了解包括光敏电阻等数种传感器元件的原理、特性。
		
		\item 学习使用 ELVIS 并了解其测量原理。
		
		\item 练习规范连接电路、测量电压与电流。
		
		\item 练习规范记录实验结果和数据处理。
	\end{enumerate}
	
	\subsection{仪器用具}
	\begin{table}[htbp]
		\centering
		\renewcommand\arraystretch{1.6}
		% \setlength{\tabcolsep}{10mm}
		\begin{tabular}{p{0.05\textwidth}|p{0.20\textwidth}|p{0.05\textwidth}|p{0.5\textwidth}}
			\hline
			编号& 仪器用具名称 & 数量 &  主要参数（型号，测量范围，测量精度等） \\
			\hline
			1& ELVIS传感器实验电路板 & 1 & -- \\
			
			2& 传感器实验套件 & 1 & 光敏电阻、PIN二极管 \\
			
			3& 鳄鱼夹表笔套 & 2 & -- \\
			
			4& 面包板测试线 & 20 & -- \\
			\hline
		\end{tabular}
	\end{table}
	
	\subsection{实验内容}
	\begin{enumerate}
		\item 光敏电阻特性实验: 定量记录结果；
		\item PIN 二极管特性实验: 定量记录结果；
		\item 霍尔 IC 电机转速实验: 定量记录结果。		
	\end{enumerate}
	
	\subsection{原理概述}
	光敏电阻是一种利用半导体的光电导效应制成的电阻器，它的电阻值能随着外界光照强弱而变化。光电导效应是指半导体材料在光的作用下，其导电能力发生变化的现象。当光子的能量大于半导体的禁带宽度时，价带中的电子吸收光子的能量后跃迁到导带，形成电子-空穴对，增加了半导体中的载流子浓度，从而降低了半导体的电阻。光照停止后，电子-空穴对逐渐复合，半导体的电阻恢复原值。
	
	\begin{figure}[htbp]
		\centering
		\includegraphics[width=0.5\textwidth]{AA4GraAD1.png}
		\caption{光敏电阻实验电路原理图}
	\end{figure}
	
	光敏电阻的基本特性包括伏安特性、光照特性、光谱特性、频率特性和温度特性。伏安特性是指在一定照度下，光敏电阻两端的电压与流过的电流之间的关系。光照特性是指在一定的外加电压下，光敏电阻的光电流与光通量之间的关系。光谱特性是指光敏电阻对不同波长的光的相对灵敏度。频率特性是指光敏电阻对脉冲光的响应速度。温度特性是指光敏电阻的电阻值随环境温度的变化而变化。	
	
	光敏电阻特性实验的目的是通过测量光敏电阻的伏安特性曲线和光照特性曲线，了解光敏电阻的工作原理和基本特性，掌握光路的调整方法，提高在光电效应原理、光敏电阻工作原理和基本特性以及应用方面的资料信息的选择搜集能力。实验仪器是一种测量光敏电阻基本特性的实验装置，包括光源、两个聚光镜、偏振器、接收器等。实验内容是在一定照度下，测量光敏电阻电压与电流的关系曲线；在一定工作电压下，测量光敏电阻的照度与电流的关系曲线。实验数据记录和处理后，可以绘制出光敏电阻的伏安特性曲线和光照特性曲线，分析其规律和特点，总结实验结果和结论。
	
	\begin{figure}[htbp]
		\centering
		\includegraphics[width=0.5\textwidth]{AA4GraAD2.png}
		\caption{光敏电阻实验仪器界面图}
	\end{figure}
	
	PIN二极管是一种光电探测器，它可以将光信号转换为电信号。它的结构是在P型半导体和N型半导体之间插入一层本征半导体（I层），形成一个PN结。I层的作用是增加耗尽区的宽度，提高光电转换的效率和灵敏度，同时降低结电容，提高频率响应。
	
	\begin{figure}[htbp]
		\centering
		\includegraphics[width=0.5\textwidth]{AA4GraAD3.png}
		\caption{PIN特性曲线}
	\end{figure}
	
	PIN二极管特性实验主要有两个部分：静态特性和动态特性。静态特性是指在不同的光照强度和偏压下，测量PIN二极管的光电流和响应度。动态特性是指在不同的光照波长和偏压下，测量PIN二极管的光谱响应和截止频率。
	
	静态特性实验的原理是：在PIN二极管的两端加上一个直流电源，调节电源的电压，使PIN二极管处于正向或反向偏置状态。然后用一个光源照射PIN二极管，调节光源的光照强度，使PIN二极管产生光电流。用一个电流表测量光电流的大小，用一个电压表测量PIN二极管两端的电压。根据光电流和电压的关系，可以计算出PIN二极管的响应度，即光电流与光功率的比值。响应度反映了PIN二极管的光电转换效率，响应度越高，效率越高。
	
	动态特性实验的原理是：在PIN二极管的两端加上一个直流电源，调节电源的电压，使PIN二极管处于反向偏置状态。然后用一个单色光源照射PIN二极管，调节光源的波长，使PIN二极管产生光电流。用一个电流表测量光电流的大小，用一个电压表测量PIN二极管两端的电压。根据光电流和波长的关系，可以绘制出PIN二极管的光谱响应曲线，即光电流与波长的函数关系。光谱响应反映了PIN二极管对不同波长的光的灵敏度，光谱响应越宽，灵敏度越高。同时，根据PIN二极管的结电容和电阻，可以计算出PIN二极管的截止频率，即PIN二极管能够正常工作的最高频率。截止频率反映了PIN二极管的频率响应能力，截止频率越高，能力越强。
	
	\begin{figure}[htbp]
		\centering
		\includegraphics[width=0.5\textwidth]{AA4GraAD4.png}
		\caption{PIN实验电路原理图}
	\end{figure}
	
	\begin{figure}[htbp]
		\centering
		\includegraphics[width=0.5\textwidth]{AA4GraAD5.png}
		\caption{PIN实验仪器界面图}
	\end{figure}
	
	霍尔 IC 是一种利用霍尔效应将磁场信号转换为电信号的集成电路，它由霍尔元件、放大器、比较器和恒流源等组成。霍尔效应是指当一个导体通过电流时，如果在垂直于电流方向施加一个磁场，那么在导体的两侧就会产生一个垂直于电流和磁场方向的电势差，这个电势差称为霍尔电压，它与磁场强度成正比。
	
	基于霍尔 IC 的直流电机转速测量实验的原理是：在电机的转轴上固定一个带有多个磁极的磁环，使之与霍尔 IC 保持一定的距离。当电机转动时，磁环上的磁极会依次经过霍尔 IC，使其感受到周期性变化的磁场。霍尔 IC 根据磁场的变化，输出相应的电压信号，经过放大和整形后，形成与磁极数目相同的方波脉冲信号。这些脉冲信号的频率与电机的转速成正比，因此，通过测量脉冲信号的频率，就可以得到电机的转速。
	
	\begin{figure}[htbp]
		\centering
		\includegraphics[width=0.5\textwidth]{AA4GraAD6.png}
		\caption{霍尔IC实验电路原理图}
	\end{figure}
	
	\subsection{实验前思考题}
	\begin{question}
		什么是暗电阻、亮电阻？什么是暗电流、亮电流、光电流？什么是偏置电压？（3分）（分别描述描述暗电阻、亮电阻、暗电流、亮电流、光电流和偏置电压）
	\end{question}
	\begin{enumerate}
		\item \textbf{暗电阻、暗电流}：光敏电阻在不受光照射时的阻值称为暗电阻，此时流过的电流称为暗电流。
		\item \textbf{亮电阻、亮电流}：光敏电阻在受光照射时的电阻称为亮电阻，此时流过的电流称为亮电流。
		\item \textbf{光电流}：亮电流与暗电流之差称为光电流。
		\item \textbf{偏置电压}：偏置电压是用于设置器件（如晶体管、二极管、传感器等）工作点或偏置点的直流电压。在光敏元件中，适当的偏置电压可以影响其在不同光照条件下的响应和灵敏度，确保它处于最佳工作状态。
		
		（参考资料：偏置电压（Bias Voltage）是在电子电路中用来设置器件工作点或偏置点的电压。它是一种直流（DC）电压，用于确保电路中的半导体器件（如晶体管、运放等）在其工作范围内正常工作。在许多电子元件中，特别是对于需要在特定工作点处线性放大信号的器件，需要提供适当的偏置电压。对于晶体管而言，例如双极型晶体管（BJT），提供适当的偏置电压能够确保其在放大信号时处于线性工作区域，从而避免失真。同样，对于场效应晶体管（FET）或运算放大器等器件，偏置电压的设置对于其性能和稳定性也至关重要。偏置电压可以是固定的或可调节的，具体取决于电路的设计需求。它可以通过外部电路提供，如电阻分压网络、电源电压、专门的偏置电路等来生成。在设计电路时，合适的偏置电压对于确保器件在工作点处能够稳定地工作非常重要。过高或过低的偏置电压都可能导致器件工作不稳定、失真或不正常的行为。因此，在电路设计过程中，需要仔细考虑和计算合适的偏置电压，以保证电路的性能和稳定性。）
	\end{enumerate}
	
	\clearpage
	\begin{question}
		画出测量光敏电阻的简单电路图。（3分）（绘制测量暗电阻、亮电阻和伏安特性的电路原理图）
	\end{question}
	如\cref{fig:fig1}、\cref{fig:fig2}、\cref{fig:fig3}。
	\begin{figure}[htbp]
		\centering
		\includegraphics[width=0.3\textwidth]{AA4Gra1.png}
		\caption{光敏电阻：测量暗电阻原理图}
		\label{fig:fig1}
	\end{figure}
		
	\begin{figure}[htbp]
		\centering
		\includegraphics[width=0.3\textwidth]{AA4Gra2.png}
		\caption{光敏电阻：测量亮电阻原理图}
		\label{fig:fig2}
	\end{figure}
	
	\begin{figure}[htbp]
		\centering
		\includegraphics[width=0.3\textwidth]{AA4Gra3.png}
		\caption{光敏电阻：测量伏安特性原理图}
		\label{fig:fig3}
	\end{figure}
	
	\clearpage
	\begin{question}
		画出测量 PIN 二极管的简单电路图。（2分）（绘制测量暗电流、亮电流的电路原理图）
	\end{question}
	如\cref{fig:fig4}、\cref{fig:fig5}所示。
	\begin{figure}[htbp]
		\centering
		\includegraphics[width=0.5\textwidth]{AA4Gra4.png}
		\caption{PIN：测量暗电流原理图}
		\label{fig:fig4}
	\end{figure}
	
	\begin{figure}[htbp]
		\centering
		\includegraphics[width=0.5\textwidth]{AA4Gra5.png}
		\caption{PIN：测量亮电流原理图}
		\label{fig:fig5}
	\end{figure}
	
	\begin{question}
		什么是霍尔效应？什么是单级霍尔效应开关？（4分）（霍尔效应2分，单极霍尔效应开关2分）
	\end{question}
	简单来说，霍尔效应是指当固体导体放置在一个个磁场内，且有电流通过时，导体内的电荷载子受到洛伦兹力而 偏向一边，继而产生电压(霍尔电压)的现象。单级霍尔效应开关是一种利用霍尔效应作为开关的器件，它通常由一个个磁性材料和一个晶体管组 成。当当磁性材料受到足够强的的磁场时，晶体管导通，输出低电平;当当磁场减弱时，晶体管截止，输出高电平。
	具体来说，\textbf{霍尔效应}是指当电流通过导体时，放置在该导体附近的磁场会产生的一种现象。当导体中的自由电荷（例如电子）受到磁场的影响时，它们会受到力的作用，导致在该导体的一侧产生电势差。（当一段导体（通常是金属或半导体）中有电流流动时，在该导体两侧会形成一个电势差，这个现象就是霍尔效应。这个电势差的大小与电流本身、外加磁场的强度和导体材料的性质相关。）这种现象是由美国物理学家爱德华·霍尔在1879年首次描述和观察到的，因此被称为霍尔效应。
	\textbf{单级霍尔效应开关}是利用霍尔效应设计的一种开关装置。它包含一个霍尔元件（霍尔传感器），可以感知周围磁场的变化并产生相应的输出信号。这种霍尔元件一般包括霍尔板、电源和输出接口。单级霍尔效应开关通常在检测磁场的应用中使用。当外部磁场达到或超过预设阈值时，霍尔传感器会感知到磁场变化，并输出一个信号，通常是电压或电流信号，以指示磁场的存在或状态变化。这种开关常用于接近传感器、位置检测、磁场控制等应用中，因为它们可以快速、准确地检测磁场的变化，并触发相应的操作或控制。
	
	\begin{question}
		简述 PID 控制器的工作原理（4分）（用自己的语言分别描述P、I、D的意义以及PID控制器的工作原理）
	\end{question}
	PID控制器是一种常用的反馈控制器，通过调节输出信号来控制系统，使其达到期望的状态或保持在期望值附近。PID代表比例（Proportional）、积分（Integral）、微分（Derivative）三个部分，其中 P 代表“比例”控制（根据误差和设定值的比例关系来调节控制量，使系统的输出趋于设定值），I 代表“积分”控制（根据误差在一段时间内的累积变化来调节控制量，使系统的输出能够消除静差或稳态误差），D 代表“微分”控制（根据误差在一段时间内的变化率来调节控制量，使系统能够预测未来的误差趋势），它们分别对系统的当前误差、历史误差和误差变化率进行处理。
	\begin{enumerate}
		\item 比例（Proportional）P：比例控制作用于当前误差，其输出与当前误差成正比。P部分决定了输出信号与偏差的大小成比例的关系。增加比例控制可以提高系统的响应速度，但可能导致系统超调或稳定性问题。
		\item 积分（Integral）I：积分控制用于历史误差的累积，它的作用是消除系统稳态误差。I部分会积累偏差的历史信息并用于调整输出信号，使系统更快地达到稳定状态。但过大的积分部分可能导致系统超调或震荡。
		\item 微分（Derivative）D：微分控制用于预测误差的变化趋势，它可以减小系统的过渡过程中的振荡和超调。D部分根据误差变化率的斜率来调整输出信号，有助于抑制系统的振荡和提高系统的稳定性。然而，过大的微分作用可能会导致系统对噪声过于敏感。
	\end{enumerate}
	PID控制器的工作原理是将这三个部分的输出信号进行加权求和，生成最终的控制信号，用于调整系统的操作。PID控制器的输出信号由当前误差、历史误差和误差变化率三部分综合决定，通过不断地调整输出信号来使系统的实际输出与期望值尽可能接近，从而实现对系统的精确控制。PID控制器是一种广泛应用于工业控制和自动化系统中的控制算法，其通过不断地调整控制信号来维持系统在设定点或期望状态附近。
	
	%实验记录	
	\clearpage
	\begin{table}
		\renewcommand\arraystretch{1.7}
		\centering
		\begin{tabularx}{\textwidth}{|X|X|X|X|}
			\hline
			专业： & 物理学 & 年级： & 2022级 \\
			\hline
			姓名： & 杨舒云 & 学号： & 22344020\\
			\hline
			室温： & 23摄氏度 & 实验地点： & 实验室507 \\
			\hline
			学生签名：& 杨舒云 & 评分： &\\
			\hline
			实验时间：& 2023/11/16 & 教师签名：&\\
			\hline
		\end{tabularx}
	\end{table}
	
	\section{AA4 传感器综合实验  \quad\heiti 实验记录}
	\subsection{实验内容、步骤与结果}
	\subsubsection{光敏电阻实验}
	由于仪器灵敏度较差、耗费时间较长，不再测量光敏器件的暗电流。仅观察其光照特性。
	\begin{enumerate}
		\item 观测光敏电阻的伏安特性（5分）（每组数据未填写-1分）
		
		按照讲义连好电路。将面板上的“电流表”插孔连线到左侧电流测量表。检查无误后，再打开ELVIS PROTO 电源。
		采用白光照射光敏电阻（应使得光源尽可能靠近电阻，白光直射电阻；如引脚过长，请折弯引脚）。调节偏置电源电压为 4V、6V、8V、10V、12V（在其附近即可），测量电路中电流的大小，记录如\cref{tab:tab1}所示。
		
		\begin{table}[h]
			\centering
			\caption{测量结果：光敏电阻的伏安特性}
			\label{tab:tab1}
			\begin{tabular}{|c|c|c|c|c|c|}
				\hline
				偏置电源电压/V (R=20kΩ) & 2 & 4 & 6 & 8 & 10 \\
				\hline
				通过传感器电流/mA (R=20kΩ) & 0.082 & 0.164 & 0.246 & 0.328 & 0.410 \\
				\hline
				偏置电源电压/V (R=1kΩ) & 4 & 6 & 8 & 10 & 12 \\
				\hline
				通过传感器电流/mA (R=1kΩ) & 0.93 & 1.38 & 1.83 & 2.31 & 2.77 \\
				\hline
			\end{tabular}
		\end{table}
		
		
		\item 测定光敏电阻的光照特性（5分）（每组数据未填写-0.5分，向上取整）
		
		请将面板上的“电压表”插孔连线到左侧电压测量表。
		请选用一种固定颜色的光源（红、绿、蓝、橙均可），\textbf{偏置电压固定为 8 V}，将占空比从 1\%-99\% 均匀调节（记录频率设定），\textbf{$RL=1k\Omega$}（或选用合适的RL），记录测量电压随光照改变的变化（\cref{tab:tab2}）。
		
		\begin{table}[h]
			\centering
			\caption{测量结果：光敏电阻的光照特性}
			\label{tab:tab2}
				\begin{tabular}{|*{12}{c|}}
					\hline
					占空比 & 99\% & 90\% & 80\% & 70\% & 60\% & 50\% & 40\% & 30\% & 20\% & 10\% & 1\% \\
					\hline
					电压/V & 0.991 & 0.925 & 0.844 & 0.761 & 0.671 & 0.579 & 0.490 & 0.396 & 0.296 & 0.182 & 0.030 \\
					\hline					
				\end{tabular}
		\end{table}
		
		\item 测定光敏电阻的光谱特性 （4分）（每组数据未填写-1分）
		
		请调节占空比 100\%，\textbf{偏置电压固定为 8V}。使用至少三种不同颜色的 LED 光源照射光敏电阻/光敏二极管，测量电压（\cref{tab:tab3}）。
		
		\begin{table}[h]
			\centering
			\caption{测量结果：光敏电阻的光谱特性}
			\label{tab:tab3}
			\begin{tabular}{|c|c|c|c|}
				\hline
				光源颜色 & 红 & 绿 & 黄 \\
				\hline
				电压1/V & 1.015 & 0.657 & 0.227 \\
				\hline
				电压2/V & 0.920 & 0.608 & 0.274 \\
				\hline
				电压3/V & 0.992 & 0.606 & 0.264 \\
				\hline
			\end{tabular}
		\end{table}
		
	\end{enumerate}
	
	附：原始数据如\cref{fig:figA1}所示。
	
	\begin{figure}[htbp]
		\centering
		\includegraphics[width=0.5\textwidth]{AA4GraA1.jpg}
		\caption{原始数据1}
		\label{fig:figA1}
	\end{figure}
	
	数据分析见\textbf{实验数据分析}部分。
	
	\clearpage
	\subsubsection{PIN 二极管特性实验}
	\begin{enumerate}
		\item 观测 PIN 光电二极管的光照特性（4分）（每组数据未填写-0.5分，向上取整）
		
		参照讲义中的测量电路，选择合适的测量方法（电压/电流），设计合适的 RL并记录\textbf{（$RL=1k\Omega$）}，\textbf{选择E=12V}。调整不同的光照占空比，测得的光生电流并填入\cref{tab:tab4}中。
		
		\begin{table}[h]
			\centering
			\caption{测量结果：PIN的光照特性}
			\label{tab:tab4}
			\begin{tabular}{|c|c|c|c|c|c|c|}
				\hline
				光照占空比 & 0 & 20\% & 40\% & 60\% & 80\% & 100\% \\
				\hline
				光生电流/μA（白） & 0.4 & 8.6 & 17.1 & 25.5 & 33.9 & 41.9 \\
				\hline
				光生电流/μA（红） & 0.5 & 8.8 & 17.2 & 26.4 & 34.2 & 40.4 \\
				\hline
				电压/mV（红） & 0.5 & 9.7 & 19.0 & 28.1 & 37.5 & 46.0 \\
				\hline
				电压/mV（白） & 0.4 & 7.9 & 16.0 & 24.1 & 31.7 & 39.3 \\
				\hline
			\end{tabular}
		\end{table}
		
		\item 观测 PIN 光电二极管的伏安特性（4分）（每组数据未填写-0.5分，向上取整）
		
		根据上题中的测量电路，RL 选择合适值并记录\textbf{（$RL=1k\Omega$）}，\textbf{占空比50\%}，保持光照度不变，分别测量计算不同 E 值时的光生电流，填入\cref{tab:tab5}。
		
		然后调整占空比为 100\%，保持光照度不变，分别测量不同 E 值时的光生电流，填入\cref{tab:tab5}。
		
		\begin{table}[h]
			\centering
			\caption{测量结果：PIN的伏安特性}
			\label{tab:tab5}
			\begin{tabular}{|c|c|c|c|c|c|c|c|}
				\hline
				偏压/V & 0 & -2 & -4 & -6 & -8 & -10 & -12 \\
				\hline
				光生电流/mA（50\%） & 0.014 & 1.330 & 3.200 & 5.100 & 7.020 & 8.940 & 10.810 \\
				\hline
				光生电流/mA（100\%） & 0.013 & 1.310 & 3.170 & 5.030 & 6.920 & 8.880 & 10.830 \\
				\hline
			\end{tabular}
		\end{table}
		
	\end{enumerate}
	附：原始数据如\cref{fig:figA2}所示。
	
	\begin{figure}[htbp]
		\centering
		\includegraphics[width=0.5\textwidth]{AA4GraA2.jpg}
		\caption{原始数据2}
		\label{fig:figA2}
	\end{figure}
	
	数据分析见\textbf{实验数据分析}部分。
	
	\clearpage
	\subsubsection{霍尔 IC 测速实验步骤与记录}
	
	\begin{enumerate}
		\item 观测开环系统直流电机转速与控制电压的关系（4分）（每组数据未填写-0.5分，向上取整）
		
		按照讲义中的测量电路，调整不同的驱动电压+VPS，测得不同的6组转速（电机不转不需要记录）填入\cref{tab:tab6}中。
		
		\begin{table}[h]
			\centering
			\caption{测量结果：霍尔IC}
			\label{tab:tab6}
			\begin{tabular}{|c|c|}
				\hline
				驱动电压/V & 电机转速 \\
				\hline
				2.70 & 2360 \\
				3.40 & 3438 \\
				3.60 & 3750 \\
				4.00 & 4331 \\
				4.20 & 4600 \\
				4.95 & 5700 \\
				\hline
			\end{tabular}
		\end{table}
		
				
		\item 观测闭环系统 PID 控制参数与转速调节的关系（4分）（配置3种控制器，并标注清楚）
		
		根据讲义中电路图连接电路，设定合适的目标转速为，分别合理配置比例+积分+微分（PID）控制器，观测不同PID参数下的阶跃运动控制曲线。	
		
		利用测得的数据作图详见\textbf{数据分析}部分，以下展示一幅图（\cref{fig:fig12}）。
		
		\begin{figure}[htbp]
			\centering
			\includegraphics[width=0.6\textwidth]{AA4Gra12.jpg}
			\caption{阶跃响应曲线}
			\label{fig:fig12}
		\end{figure}
	\end{enumerate}
	
	附：原始数据如\cref{fig:figA3}所示。
	
	\begin{figure}[htbp]
		\centering
		\includegraphics[width=0.9\textwidth]{AA4GraA3.jpg}
		\caption{原始数据3}
		\label{fig:figA3}
	\end{figure}
	
	数据分析见\textbf{实验数据分析}部分。
	
	
	\subsection{原始数据记录}
	实验记录本上的原始数据见\cref{fig:figA1}、\cref{fig:figA2}和\cref{fig:figA3}（签字）。
	
	实验台桌面整理见\textbf{附件}部分(\cref{fig:figA4})。
	
	\subsection{实验过程中遇到的问题记录}
	\begin{enumerate}
		\item 实验仪器中有大量面包板测试线是损坏的但没有被清理出来，这大大影响实验者接线的进程。
		\item 观测光敏电阻光照特性时多次出现示数突变影响实验进程。
		\item 有些PID比例无法使电机转起来。
	\end{enumerate}
	
	
	
	%分析与讨论	
	\clearpage
	\begin{table}
		\renewcommand\arraystretch{1.7}
		\begin{tabularx}{\textwidth}{|X|X|X|X|}
			\hline
			专业：& 物理学 &年级：& 2022级\\
			\hline
			姓名： & 杨舒云 & 学号：& 22344020\\
			\hline
			日期：& 2023/10/19 & 评分： &\\
			\hline
		\end{tabularx}
	\end{table}
	
	\section{AA4 传感器综合实验 \quad\heiti 分析与讨论}
	\subsection{实验数据分析}
	\subsubsection{观测光敏电阻的伏安特性（3分）}
	\begin{enumerate}
		\item 通过五次测量，第一组实验数据计算得到的电阻结果如表\cref{tab:tab7}所示。
		
		\begin{table}[h]
			\centering
			\caption{电阻1：光敏电阻的伏安特性}
			\label{tab:tab7}
			\begin{tabular}{|c|c|c|c|c|c|c|c|} 
				\hline 
				偏置电源电压/V（R=20kΩ） & 2 & 4 & 6 & 8 & 10 & 均值 & 均值标准差 \\ 
				\hline 
				通过传感器电流/mA（R=20kΩ） & 0.082 & 0.164 & 0.246 & 0.328 & 0.410 & & \\
				\hline 
				电阻/kΩ & 24.390 & 24.390 & 24.390 & 24.390 & 24.390 & 24.390 & 0.000 \\ 
				\hline 
			\end{tabular}
		\end{table}
				
		得到传感器电阻平均结果为$\bar{R_1}=24.39k\Omega$。
		
		通过五次测量，第二组实验数据计算得到的电阻结果如\cref{tab:tab8}所示。
		
		\begin{table}[h]
			\centering
			\caption{电阻2：光敏电阻的伏安特性}
			\label{tab:tab8}
			\begin{tabular}{|c|c|c|c|c|c|c|c|} 
				\hline 
				偏置电源电压/V（R=1kΩ） & 4 & 6 & 8 & 10 & 12 & 均值 & 均值标准差 \\ 
				\hline 
				通过传感器电流/mA（R=1kΩ） & 0.93 & 1.38 & 1.83 & 2.31 & 2.77 & & \\ 
				\hline 
				电阻/kΩ & 4.301 & 4.348 & 4.372 & 4.329 & 4.332 & 4.336 & 0.012 \\ 
				\hline 
			\end{tabular} 
		\end{table}
		
		得到传感器电阻平均结果为$\bar{R_2}=4.336k\Omega$。 	
		
		\item 参照《误差理论与数据处理》对五次测量的结果进行误差分析，利用罗曼诺夫斯基准则，取显著度α=0.01，对数据列进行粗大误差判断（给出过程），剔除粗大误差后，给出最终算数平均值和平均值标准差。列举不少于三项测量过程中的系统误差来源。
		
		根据罗曼诺夫准则，我们没有任何理由怀疑第一组数据中的任何一个数据是粗大误差。
		
		而对于第二组，认为8V所对应的数据有粗大误差，剔除后重新计算：均值$\bar{R'}=4.328k\Omega$，测量列标准差$\sigma'=0.019k\Omega$，取显著度$\alpha = 0.01$， 测量次数$n = 5$， 查表得$K(5, 0.01) = 6.53$，由于（数值上）
		\[|\bar{R}-\bar{R'}|=0.0088<K\sigma'\]
		
		则无充分理由认为该数据存在粗大误差，应当予以保留。
		
		观察其他数据，均在平均值附近，不应认为有粗大误差。
		
		\item 综上所述，第一组数据的平均值$\bar{R_1}=24.39k\Omega$，均值标准差为$\sigma'=0k\Omega$；第二组数据的平均值$\bar{R_2}=4.336k\Omega$，均值标准差为$\sigma'=0.012k\Omega$。
		
		\item 系统误差来源：
		\begin{itemize}
			\item 实验采用电流表测电流，所测得的电流不是真实值，而会有电流表带来的系统误差（高中知识伏安法）；
			\item 光敏电阻的表面积和光源的距离会影响光敏电阻的阻值，从而影响伏安特性曲线的形状；
			\item 光敏电阻的响应时间和恢复时间会导致光敏电阻的阻值不能及时跟随光强的变化，从而影响伏安特性曲线的准确性；
			\item 光敏电阻的温度特性会使光敏电阻的阻值随温度的变化而变化，从而影响伏安特性曲线的稳定性；
			\item 电压表和电流表的内阻和量程会影响测量的电压和电流的准确度，从而影响伏安特性曲线的精度；
			\item 电源的稳定性和波动性会影响测量的电压和电流的一致性，从而影响伏安特性曲线的重复性。
		\end{itemize}
	\end{enumerate}
	
	\clearpage
	\subsubsection{测定光敏电阻的光照特性（6分）}
	\begin{enumerate}
		\item 按照电路图，光敏电阻电阻 R 与电压表读数 U 的关系为（列出公式）（3分）
		\[R=\dfrac{E-U}{U}R_L\]
		
		计算得到如\cref{tab:tab9}所示结果。
		
		\begin{table}[h]
			\centering
			\caption{R 随占空比变化关系}
			\label{tab:tab9}
			\begin{tabular}{|c|c|c|c|c|c|c|c|c|c|c|c|} 
				\hline 
				占空比 & 99\% & 90\% & 80\% & 70\% & 60\% & 50\% & 40\% & 30\% & 20\% & 10\% & 1\% \\ 
				\hline 
				电压/V & 0.991 & 0.925 & 0.844 & 0.761 & 0.671 & 0.579 & 0.490 & 0.396 & 0.296 & 0.182 & 0.030 \\
				\hline 
				电阻/kΩ & 7.073 & 7.649 & 8.479 & 9.512 & 10.923 & 12.817 & 15.327 & 19.202 & 26.027 & 42.956 & 265.667 \\ 
				\hline 
			\end{tabular}
		\end{table}
		
		\item 由此计算得到 R 随占空比变化的曲线（作图）
		
		如\cref{fig:fig6}所示。
		
		\begin{figure}[htbp]
			\centering
			\includegraphics[width=0.6\textwidth]{AA4Gra6.jpg}
			\caption{R 随占空比变化的曲线}
			\label{fig:fig6}
		\end{figure}
	\end{enumerate}
	
	\clearpage
	\subsubsection{测定光敏电阻的光谱特性 （3分）（每组数据未填写-0.5分，向上取整，变化图1分）}
	如\cref{tab:tab10}所示。
	
	\begin{table}[h]
		\centering
		\caption{R 随占空比变化关系}
		\label{tab:tab10}
		\begin{tabular}{|c|c|c|c|} 
			\hline 
			光源颜色 & 红 & 绿 & 黄 \\ 
			\hline 
			电压1/V & 1.015 & 0.657 & 0.227 \\ 
			\hline 
			电压2/V & 0.920 & 0.608 & 0.274 \\
			\hline 
			电压3/V & 0.992 & 0.606 & 0.264 \\ 
			\hline 
			电压均值/V & 0.976 & 0.624 & 0.255 \\ 
			\hline 
			波长/nm & 635.000 & 510.000 & 465.000 \\ 
			\hline 
			电阻/kΩ & 7.199521695 & 11.82736505 & 30.37254902 \\ 
			\hline 
		\end{tabular}
	\end{table}	
	
	由此得到电阻随光源波长变化图如\cref{fig:fig7}所示。
	
	\begin{figure}[htbp]
		\centering
		\includegraphics[width=0.6\textwidth]{AA4Gra7.jpg}
		\caption{电阻随光源波长变化图}
		\label{fig:fig7}
	\end{figure}
	
	\subsubsection{作图并分析 PIN 光电二极管的光照特性曲线。（5分）（作图2分，分析3分）}
	\begin{figure}[htbp]
		\centering
		\subfloat[]{
			\includegraphics[width=0.5\textwidth]{AA4Gra8.jpg}
		}
		\subfloat[]{
			\includegraphics[width=0.5\textwidth]{AA4Gra9.jpg}
		}
		\caption{PIN光照特性曲线}
		\label{fig:fig8}			
	\end{figure}
	
	\begin{figure}[htbp]
		\centering
		\subfloat[]{
			\includegraphics[width=0.5\textwidth]{AA4Gra10.jpg}
		}
		\subfloat[]{
			\includegraphics[width=0.5\textwidth]{AA4Gra11.jpg}
		}
		\caption{PIN光照特性曲线}
		\label{fig:fig10}			
	\end{figure}
	
	分别对红光、白光进行了电流、电压共四组测量，作图并进行线性拟合后如\cref{fig:fig8}、\cref{fig:fig10}所示，可以发现它们都有着良好的线性关系，光生电流与光照占空比成线性正相关，光照占空比越大，光生电流越强。
	
	\clearpage
	\subsubsection{根据测得的实验数据，作出三组PID 控制器下的阶跃响应曲线， 并进行分析比较。（7分）（作出3组曲线，每组2分，分析比较1分）}
	根据测得的数据，作图如\cref{fig:fig15}、\cref{fig:fig13}和\cref{fig:fig14}所示。
	
	\begin{figure}[htbp]
		\centering
		\includegraphics[width=0.9\textwidth]{AA4Gra15.jpg}
		\caption{阶跃响应曲线：P=0.00075；I=0.0083；D=0.0011}
		\label{fig:fig15}
	\end{figure}
	
	\begin{figure}[htbp]
		\centering
		\includegraphics[width=0.7\textwidth]{AA4Gra13.jpg}
		\caption{阶跃响应曲线：P=0.001；I=0.0058；D=0.0063}
		\label{fig:fig13}
	\end{figure}
	
	\begin{figure}[htbp]
		\centering
		\includegraphics[width=0.7\textwidth]{AA4Gra14.jpg}
		\caption{阶跃响应曲线：P=0.0007；I=0.0058；D=0.0035}
		\label{fig:fig14}
	\end{figure}
	
	\clearpage
	阶跃运动控制曲线是系统在控制器输出阶跃变化作用下，过程变量从初始状态到稳定状态相对于时间的曲线。阶跃运动控制曲线可以反映系统的动态响应和稳态误差，是评价系统性能的重要指标。
	
	\begin{enumerate}
		\item 第一幅图（\cref{fig:fig15}）P=0.00075、I=0.0083、D=0.0011，第二幅图（\cref{fig:fig13}）P=0.001、I=0.0058、D=0.0063，第三幅图（\cref{fig:fig14}）P=0.0007、I=0.0058、D=0.0035；
		
		\item 这三种情况PID大致配比（多少）相同，因此所得到的跃阶曲线是同一类型的，即系统的稳态误差较小，系统的超调量接近零，系统的最终转速与目标转速趋向一致，具有不错的精度；
		
		\item 差别会体现在调节时间上，第一种情况系统的调节时间约在20秒以上，具有一定的快速性；第二种情况系统的调节时间约在30秒以上，并且此后出现了反复的波动，快速性和调节性相对较差；第三种情况系统的调节时间约在15秒以上，且后续稳定，具有良好的快速性。
		
		\item 进一步对比各图，第二种情况有着较多较大振荡幅度，这说明 PID 控制器的比例系数过大，积分时间过长，微分时间过短，导致系统不稳定，无法快速收敛到设定值。第三种情况的响应曲线振荡幅度较小，振荡数较少，这说明 PID 控制器的比例系数适中，积分时间适中，微分时间适中，使得系统具有较好的动态性能和稳态性能。第一种情况介于二者之间。
			
	\end{enumerate}
	
	
	
	\clearpage
	\subsection{实验后思考题}
	\begin{question}
		试举光敏电阻在日常生活中的应用事例。（3分）
	\end{question}
	光敏电阻是一种根据光的强度改变其电阻值的器件，它在日常生活中有很多应用。
	
	光敏电阻可以用作光敏传感器，用于检测光的存在和强度。例如，一些自动开关灯的装置就是利用光敏电阻来控制电路的开关，当光线变暗时，光敏电阻的电阻增大，导致电路通电，灯就会亮起；当光线变亮时，光敏电阻的电阻减小，导致电路断电，灯就会熄灭。
	
	光敏电阻也可以用于测量光的强度，例如，一些照相机的测光表就是利用光敏电阻来测量光线的亮度，从而调节快门速度和光圈大小，以达到最佳的曝光效果。
	
	光敏电阻的滞后特性可以用于音频压缩和外部感应，例如，一些电子琴的音量控制就是利用光敏电阻来实现的，当光线变暗时，光敏电阻的电阻增大，导致电子琴的音量减小；当光线变亮时，光敏电阻的电阻减小，导致电子琴的音量增大¹。
	
	光敏电阻还可以用于一些与光相关的设备，
	例如，闹钟，户外时钟，太阳能路灯等，它们都是利用光敏电阻来感应光线的变化，从而实现不同的功能。
	
	\begin{question}
		调研回答，光敏电阻的阻值随光照迅速减小的物理机制？（3分）
	\end{question}
	\begin{itemize}
		\item 光敏电阻的工作原理是利用了光电导效应，即当光子照射到半导体材料上时，会激发出电子-空穴对，从而增加材料的导电性。光敏电阻的阻值与光照强度成反比，即光照越强，阻值越小。
		
		\item 光敏电阻的材料可以分为本征型和掺杂型两种。本征型光敏电阻使用未掺杂的硅或锗等材料，光子需要有足够的能量才能将价带上的电子激发到导带上，从而产生自由电子和空穴。这种类型的光敏电阻对可见光和紫外光比较敏感。掺杂型光敏电阻使用掺入杂质的材料，如硫化镉、硒化铅等，杂质会在价带上形成一个新的能级，这个能级上的电子只需要较小的能量就能跃迁到导带上，从而增加导电性。这种类型的光敏电阻对红外光比较敏感。
		
		\item 不同类型的光敏电阻对不同波长的光有不同的灵敏度，如果光的波长超出了一定的范围，就不会对光敏电阻的阻值产生影响。因此，光敏电阻的光谱响应曲线是反映了光的波长与光敏电阻的灵敏度之间的关系。		
	\end{itemize}
	
	
	
	\clearpage
	\section{AA4 传感器综合实验 \quad\heiti 参考文献与附件}
	\subsection{其他（20分）}
	1、作业整洁2分；
	
	2、实验后整理桌面3分；
	
	3、实验报告信息填写1分；
	
	4、现场实验表现14分：注意在现场实验中，经老师二次以上提示未能顺利完成实验的扣3分，违规操作造成仪器损坏的扣8分，违规操作造成安全事故的得0分。
	
	\subsection{参考文献}
	[1]维基百科
	
	[2]实验讲义
	
	\subsection{附件}
	\begin{figure}[htbp]
		\centering
		\includegraphics[width=0.7\textwidth]{AA4GraA4.jpg}
		\caption{实验台桌面整理}
		\label{fig:figA4}
	\end{figure}
	
\end{document}